\centerline{\bf \Huge  Подсчёт двумя способами}
\title{Подсчёт двумя способами}

\begin{flushright}
    \scriptsize Если я сказал — не брал, значит — не отдам!\\Урахара Киске. Bleach
\end{flushright}

\problem{0.1} Можно ли в прямоугольной таблице 5×10 так расставить числа, чтобы сумма чисел каждой строки равнялась бы 30, а сумма чисел каждого столбца равнялась бы 10?

\problem{0.2} Три купчихи – Сосипатра Титовна, Олимпиада Карповна и Поликсена Уваровна – сели пить чай. Олимпиада Карповна и Сосипатра Титовна выпили вдвоём 11 чашек, Поликсена Уваровна и Олимпиада Карповна – 15, а Сосипатра Титовна и Поликсена Уваровна – 14. Сколько чашек чая выпили все три купчихи вместе?

\problem{0.3} Несколько шестиклассников и семиклассников обменялись рукопожатиями. При этом оказалось, что каждый шестиклассник пожал руку семи семиклассникам, а каждый семиклассник пожал руку шести шестиклассникам. Кого было больше - шестиклассников или семиклассников? 

\problem{1} Ковровая дорожка покрывает лестницу из 9 ступенек. Длина и высота лестницы равны 2 метрам. Хватит ли этой ковровой дорожки, чтобы покрыть лестницу из 10 ступенек длиной и высотой 2 метра?

\problem{2} Рита, Люба и Варя решали задачи. Чтобы дело шло быстрее, они купили конфет и условились, что за каждую решённую задачу девочка, решившая её первой, получает четыре конфеты, решившая второй — две, а решившая последней — одну. Девочки говорят, что каждая из них решила все задачи и получила 20 конфет, причём одновременных решений не было. Они ошибаются. Как вы думаете, почему? 

\problem{3} На сторонах шестиугольника было записано шесть чисел, а в каждой вершине – число, равное сумме двух чисел на смежных с ней сторонах. Затем все числа на сторонах и одно число в вершине стерли. Можно ли восстановить число, стоявшее в вершине? 

\problem{4} Можно ли в клетки квадрата 10×10 поставить некоторое количество звёздочек так, чтобы в каждом квадрате 2×2 было ровно две звёздочки, а в каждом прямоугольнике 3×1 – ровно одна звёздочка? (В каждой клетке может стоять не более одной звёздочки.)

\problem{5} Когда встречаются два жителя Цветочного города, один отдает другому монету в 10 копеек, а тот ему - 2 монеты по 5 копеек. Могло ли случиться так, что за день каждый из 1990 жителей города отдал ровно 10 монет?